\chapter{Introduction}

Purpose of the course \emph{Data structures and algorithms}
is to teach methods by which we can solve
\emph{effectively} computational problems.
Basic programming courses are focused
for learning programming skills.
Now it's time to go one step further and start pinning
also pay attention to how fast the algorithms work.

The efficiency of algorithms is of great importance in practice.
For example, an online route guide is useful because
that it gives a suggestion of a route right after we are
have announced where we want to travel to.
If you have to wait a minute or an hour for a suggested route,
this would greatly limit the use of the service.

In order for the route guide to work effectively, it is based on
well designed algorithm.
In this course we will learn how we can create similar algorithms ourselves.
In the course, we get to know both the theory of algorithm design and
into practice -- we want to deeply understand what algorithms are about,
but also know how to implement them in practice.

\section{What are algorithms?}

\index{algorithms}
\index{input}
\index{printout}

\emph{Algorithm} is a procedure that we can follow to solve
some computational problem.
The algorithm is given \emph{input} (\emph{input}),
which describes the case of the problem to be solved,
and the algorithm must produce \emph{output},
which is a response to the input given to it.

As an example, consider the problem,
where the input is a list containing $n$ integers and
the task is to calculate the sum of the numbers.
For example, if the input is $[2,4,1,8]$,
the desired output is 15 because $2+4+1+8=15$.
We can solve this problem with an algorithm,
which loops through the numbers and counts them
to the amount variable.

There are several possibilities for representing the operation of the algorithm.
One way is to verbally explain how the algorithm works,
as we just did.
Another way is to give in some programming language
written code that implements the algorithm.
For example, we can make the sum of the numbers count
this is how the algorithm works in Java and Python:

\begin{code}
    int sum = 0;
    for (int i = 0; i < n; i++) {
            sum+= numbers[i];
        }
    System.out.println(sum);
\end{code}

\begin{code}
    sum= 0
    for i in range(n):
    sum+= numbers[i]
    print(sum )
\end{code}

\index{pseudocode}

We can also represent the algorithm using \textbf{pseudocode} instead of an actual programming language.
This means we write code that closely resembles real programming languages but allows us to decide the exact syntax and take some liberties, making it easier to describe the algorithm.
For example, we can present the recent algorithm in pseudocode as follows:

\begin{code}
    sum = 0
    for i = 0 to n-1
    sum += numbers[i]
    print(sum)
\end{code}

There are many ways to notate pseudocode, and
we use quite a lot of Python-like pseudocode in this book.
We present algorithms using pseudocode,
because the course deals with general programming theory,
which is not related to a specific programming language.
In addition to this, we will get to know how certain algorithms
and data structures are implemented in Java and Python
in standard libraries.

\section{The basic building blocks of programming}

An intriguing aspect of programming is that even complex algorithms are built from simple components.
Next, we will go through the fundamental building blocks of programming, which form the foundation for designing algorithms.

\subsubsection{Varia}

\subsubsection{Variable}

\index{variable}
\emph{Variable} (\emph{muuttuja}) contains the data processed by the algorithm.
In pseudocode, it is common practice to use variables directly without separate declarations:


\begin{code}
    a = 5
    b = 7
    c = a+b
\end{code}

\subsubsection{Conditional Statement}

\index{conditional statement}
\emph{Conditional statement} (\emph{ehtolause}) allows the program's execution to branch.
In pseudocode, we use the following syntax:


\begin{code}
    if x%2 == 0
    print("even")
    else
    print("odd")
\end{code}

\subsubsection{Loop}

\index{loop}
\emph{Loop} repeats the code inside it.
In pseudocode, we denote a for-loop that iterates through the numbers from 1 to 100 as follows:


\begin{code}
    for i = 1 to 100
    print(i)
\end{code}

The following for-loop, in turn, iterates through the elements in a list:


\begin{code}
    for x in list
    print(x)
\end{code}

Another type of loop is the while-loop, which repeats the code as long as the condition given at the beginning of the loop holds.
For example, the following loop continues as long as the value of $x$ is at least 1:


\begin{code}
    while x >= 1
    print(x)
    x /= 2
\end{code}

\subsubsection{Array}

\index{array}
\index{index}

A fundamental data structure in programming is the \emph{array} (\emph{taulukko}),
which contains a collection of sequential elements.
We can refer to array elements using the \texttt{[]} syntax:


\begin{code}
    numbers[0] = 4
    numbers[3] = 2
\end{code}

We follow the common convention in this book where the elements of an array are indexed
as $0,1,\dots,n-1$ when the array contains $n$ elements.

An array has the properties that it contains a fixed number of elements
and allows efficient access to any element based on its index.
An array corresponds to the way data is stored in a computer's memory,
and we can use arrays to implement any other data structure.

\index{list}

The table has had a strong position in traditional programming languages, but the situation is changing.
An array is a basic data structure in Java, but in Python
instead of a table, the basic data structure is \emph{list}.
However, we can use a Python list like a table,
because the only essential difference is that its size can change.
The internal implementation of the list is based on a table,
and in Chapter 4 we will see how this happens in practice.

\subsubsection{Subroutines}

\index{subroutine}
\index{procedure}
\index{function}

In this book, we use the term \emph{procedure} (\emph{proseduuri})
for a subroutine that does not return a value,
and the term \emph{function} (\emph{funktio})
for a subroutine that returns a value.

For example, the following procedure prints the numbers $1 \dots n$:

\begin{code}
    procedure print(n)
    for i = 1 to n
    print(i)
\end{code}


The following function, on the other hand, calculates $1 \dots n$ of numbers amount:

\begin{code}
    function sum(n)
    s = 0
    for i = 1 to n
    s += i
    return s
\end{code}

The terminology related to subroutines varies depending on the programming language:
for example, in Java, the term \emph{method} is used in both cases,
while in Python, the terms \emph{function} or \emph{method} are used depending on the context.

\subsubsection{Recursion}

\index{recursion}

A common programming technique in algorithm design is
\emph{recursion} (\emph{rekursio}),
which means that a subroutine calls itself.

Here is an example of using recursion:

\begin{code}
    procedure test(n)
    if n == 0
    return
    else
    print("hello")
    test(n-1)
\end{code}


This procedure prints the line "hello" $n$ times.
For example, the call \texttt{test}(3) produces the following output:

\begin{code}
    hello
    hello
    hello
\end{code}

The idea is that if $n=0$, the procedure does nothing,
because there is nothing to print.
Otherwise, the procedure prints the line "hello"
and then calls itself with the parameter $n-1$.

We use recursion often during the course,
and little by little we get to know its possibilities in more detail.

\subsubsection{***}

We have now gone through the ingredients,
which allow us to implement \emph{any} algorithm.
It is reassuring to know that even such a small number of techniques
is sufficient for designing algorithms.
Now it's all about how we know how to \emph{apply}
these techniques in different situations.
